\documentclass[
  aspectratio=169,
]{beamer}

% Packages
\usepackage{booktabs}      % For formal tables
\usepackage{amsmath}       % AMS style math environment (e.g. \align, \gather, \cases)
\usepackage{amssymb}       % AMS style math symbols (e.g. \mathbb, \leqslant)
\usepackage{caption}       % To toggle display of captions
\usepackage{luatexja}      % For Japanese

% Bibliography
\usepackage[
  backend=biber,
  style=authoryear,  % or alphabetic, ...
  sorting=none,      % or ...
]{biblatex}

\addbibresource{references.bib}

% Macros
\newcommand{\info}[1]{\textcolor{ForestGreen}{#1}}  % To emphasize secondaly to \alert (e.g. \info{...})

% Caption
% \captionsetup[figure]{labelformat=empty}
% \captionsetup[table]{labelformat=empty}

% Theorem environments
% \setbeamertemplate{theorems}[numbered]  % Numbered theorems, lemmas, etc.

% Beamer settings
\title{2.2.1 Vertex Cover Kernelization}
\subtitle{Parameterized Algorithms pp.17--22}
\author{岡山 大輝 (\textit{ad25r010@guh.u-hyogo.ac.jp})}
\institute{兵庫県立大学大学院 情報科学研究科 東川研究室}
\date{\today}

% \AtBeginSection[]
% {
%   \begin{frame}{目次}
%     \tableofcontents[currentsection]
%   \end{frame}
% }

% Theme
\usetheme{metropolis}
\usecolortheme{default}
\metroset{block=fill}
\metroset{sectionpage=none}

% Fonts
\setbeamerfont{title}{size=\LARGE,series=\bfseries}
\setbeamerfont{frametitle}{size=\large,series=\bfseries}
\setbeamerfont{theorem}{shape=\upshape}
\setbeamerfont{date}{size=\small}
\renewcommand{\kanjifamilydefault}{\gtdefault}
\renewcommand{\familydefault}{\sfdefault}
\usefonttheme{professionalfonts}
\usefonttheme{structurebold}
\usepackage{mlmodern}
\usepackage[T1]{fontenc}

% ==============================================================================================
% ==============================================================================================

\begin{document}

\begin{frame}
	\titlepage{}
\end{frame}

\begin{frame}{What is Kernelization?}
\end{frame}

\begin{frame}{Why \textsc{Vertex Cover}?}
	\begin{itemize}
		\item
	\end{itemize}
\end{frame}

\begin{frame}{Kernelization for the \textsc{Vertex Cover} Problem}
	\textit{The Vertex Cover problem is often referred to as the Drosophila of parameterized complexity.\cite{fellows2018known}}

	\begin{itemize}
		\item \(\mathcal{O}(k^2)\) --- \cite{buss1993nondeterminism}
		\item \(3k\) --- \cite{chor2004linear}
		\item \(2k\) --- \cite{chen2001vertex}
		\item \(2k\) --- \cite{dehne2004greedy}
		\item \(2k - c\) --- \cite{soleimanfallah2011kernel}
		\item \(2k - c \log k\) --- \cite{lampis2011kernel}
		\item \(2k - \mathcal{O}(\log k)\) --- \cite{narayanaswamy2012lp}
	\end{itemize}

	\note{
		\begin{itemize}
			\item \cite{chor2004linear}は副次的な結果と思われる．
		\end{itemize}
	}
\end{frame}

\begin{frame}[allowframebreaks]{References}
	\printbibliography[heading=none]{}
\end{frame}

\end{document}
